\subsection{\textbf{Sprint 1}}

Nesta sprint, enfrentamos o desafio de não ter as \textit{stakeholders} presentes na apresentação dos \textit{mockups}. Para contornar isso, os \ac{ages} IV do nosso time decidiram organizar uma reunião extraclasse, garantindo assim que estávamos no caminho correto. Embora não tenha conseguido participar dessa reunião, o \textit{feedback} recebido foi bastante positivo.

Durante o planejamento, dividimos o time em duas grandes equipes: \textit{backend} e \textit{frontend}. Essas, por sua vez, se subdividiram em duplas ou trios para atacar as \ac{us} da sprint. Integrei a equipe de \textit{front-end}, assumindo a responsabilidade pela tela inicial do RH. Dentro dessa tarefa, encarreguei-me especificamente do carrossel da tela e, posteriormente, colaborei com o \ac{ages} I responsável por outro componente dessa mesma tela, Andrei Rech.

Ao desenvolver o carrossel, optei inicialmente por utilizar uma biblioteca que já oferecia as funcionalidades necessárias, necessitando apenas de estilização adicional. Contudo, após finalizar essa etapa e discutir com os \ac{ages} III, decidimos refatorar o código para eliminar o uso da biblioteca, favorecendo um desenvolvimento completamente nativo. Felizmente, a finalização desse componente foi rápida, graças ao meu prévio estudo das tecnologias envolvidas.

Para compor a tela inicial, eu e o \ac{ages} I, Andrei Rech, nos reunimos durante uma aula para realizar os ajustes necessários, visando alcançar a responsividade mínima e alinhar o design com os \textit{mockups} no Figma. Um ponto pendente, no entanto, foi a barra de navegação, a cargo de outra equipe. Apesar disso, deixamos nossa tela praticamente pronta, aguardando apenas a integração com o \textit{backend} e a implementação da barra de navegação.

A falta de uma apresentação direta para as \textit{stakeholders} foi um contratempo, mas o progresso alcançado e a qualidade do trabalho entregue pelo time, especialmente na primeira sprint de desenvolvimento, me deixaram particularmente satisfeito.